\scanpage{336}
\mspage{75}
% archive sheet number: 335
% ms. 1–7
Si $a$ et $b$ \inserted{dans $X$} sont reliés par une flèche de $X\Sigma^{-1}$, i.e. par un $\Sigma$-chemin, alors ils sont dans une même comp. connexe de $X$. Mais l'inverse est \unsure{égalemt} vrai, comme on voit p. ex. en prenant $\Sigma=\mathrm{Fl}(X)$, et en disant que \deleted{\unsure{les composantes connexes du}}\inserted{dans le} \uline{groupoïde} $\Pi(X)$, \deleted{\unsure{si}} deux sommets sont dans une \unsure{m} comp.
% ms. 8–10
ssi $\exists\ x\longrightarrow y$ (relation qui est non seulement transitive et \unsure{reflexive|réflexive}, mais aussi \unsure{symmétrique|symétrique} \deleted{puisque $\Pi(X)$ est}\inserted{dans} un groupoïde). Un chemin
% ms. 11
\[a\xrightarrow{c}b\]
% ms. 12–18
sera appelé aussi une \uline{homotopie} de $a$ vers $b$. On dit que $a$ et $b$ sont \uline{homotopes} s'il existe une \deleted{telle} homotopie \inserted{$c$} de $a$ vers $b$, auquel cas il existe aussi une homotopie \deleted{\unsure{``opp}} ``inverse'' $c^{\unsure{\circ|o}}$ de $b$ vers $a$. Cela signifie aussi que $a$ et $b$ sont dans la $=$ composante connexe de $X$, i.e. \unsure{définissent} le $=$ élément de $\pi_0(X)$.
% ms. 19–26
L'homotopie sera dite \uline{simple} si \unsure{$c$|elle} est de type $\Delta_+$ ou $\Delta_-$, i.e. si elle est donnée par une flèche $u:a\longrightarrow b$ ou $v:b\longrightarrow a$. Si on cherche une homotopie entre $a$ et $b$, on essaiera de la trouver la plus ``courte'' possible. Si on a une homotopie $c$ de longueur $\ell$ de $a$ vers $b$, alors en remplaçant $\tau$ par $\tau_{\mathrm{red}}$, donc $I_\tau$ par $I_{\tau_{\mathrm{red}}}$, la \unsure{sous-échelle|sous-chenille} réduite
